\chapter{Estado del Arte}

En este capítulo se presenta una revisión exhaustiva de los conceptos teóricos y tecnológicos que fundamentan este trabajo. La búsqueda y rescate (SAR, por sus siglas en inglés) mediante vehículos aéreos no tripulados (UAVs) ha evolucionado desde la simple ejecución de patrones de vuelo fijos hacia sistemas inteligentes guiados por la probabilidad y la inteligencia artificial.

\section{Teoría de Búsqueda Bayesiana}

La base matemática para la planificación de misiones de búsqueda tiene sus raíces teóricas en los trabajos desarrollados durante la Segunda Guerra Mundial, documentados posteriormente por autores como Koopman (1946) \cite{koopman}. El problema fundamental no consiste únicamente en mover un dron por el espacio, sino en decidir cómo moverlo para maximizar las posibilidades de encontrar el objetivo antes de que se agote el tiempo o la batería.

Para formalizar este problema, la teoría de búsqueda clásica define varias métricas clave que los algoritmos intentan optimizar:
\begin{itemize}
    \item \textbf{Probabilidad de Área (POA):} Es la probabilidad inicial de que la persona perdida se encuentre en una zona concreta del mapa, basándose en el terreno y su perfil.
    \item \textbf{Probabilidad de Detección (POD):} Es la probabilidad de que la cámara o el sensor del dron detecte a la persona, asumiendo que esta se encuentra justo en la zona que se está sobrevolando. Depende de factores como la altura de vuelo o la densidad de los árboles.
    \item \textbf{Probabilidad de Éxito (POS):} Es la métrica definitiva a maximizar. Se calcula multiplicando la POA por la POD:
    \begin{equation}
        POS = POA \times POD
    \end{equation}
    Un algoritmo inteligente puede decidir buscar en una zona con menos probabilidad de presencia (baja POA) si el terreno está muy despejado y es fácil ver el suelo (alta POD), garantizando así un mayor éxito global.
\end{itemize}

Durante la búsqueda, a medida que los drones sobrevuelan zonas sin encontrar a la víctima, el sistema debe actualizar la creencia sobre dónde podría estar. Este proceso de descarte y actualización continua de las probabilidades se modela mediante técnicas bayesianas, como el Filtro Recursivo Bayesiano, ampliamente utilizado en la literatura moderna de robótica SAR y en el marco de búsqueda en tiempo mínimo (\textit{Minimum Time Search} o MTS) \cite{MTS_c1, lanillos-bayesian, carabaza_ingles, tesis_carabaza}.

\section{Planificación de Rutas para UAVs en SAR}

La planificación de rutas para drones en operaciones SAR se divide principalmente en dos enfoques: Coverage Path Planning (CPP), que busca cubrir sistemáticamente un área definida, e Informative Path Planning (IPP), que maximiza la adquisición de información relevante por unidad de tiempo.

\subsection{Coverage Path Planning (CPP)}

El Coverage Path Planning tiene como objetivo generar una ruta que asegure que la huella del sensor pase sobre la totalidad de un área definida al menos una vez \cite{galceran2013}. En el contexto SAR, los equipos de rescate han empleado históricamente patrones de cobertura como espirales expansivas o barridos paralelos \cite{foo2012, hart2023}.

El método más elemental de CPP es el patrón boustrophedon o ``lawnmower'' (cortacésped), que es simple de implementar manual o autónomamente \cite{bahnemann2019}. Sin embargo, una limitación principal de los métodos de cobertura exhaustiva es su ineficiencia inherente en términos de tiempo y consumo energético, particularmente cuando se aplican a áreas grandes y complejas \cite{choset2001}.

Para mitigar estas ineficiencias, se han propuesto varias adaptaciones: descomposición del área de búsqueda en celdas más pequeñas para cobertura paralela por múltiples UAVs \cite{nielsen2019, dong2020}, y el empleo de equipos jerárquicos donde un UAV de alta altitud realiza vigilancia de amplia área mientras un UAV secundario de baja altitud realiza búsquedas detalladas en regiones de interés específicas \cite{hacho2018}.

\subsection{Informative Path Planning (IPP)}

En contraste con la cobertura exhaustiva, el Informative Path Planning busca maximizar la adquisición de información relevante por unidad de tiempo o distancia recorrida \cite{rossello2022, bai2021}. Esto se logra incorporando un marco probabilístico, como el mapa de Probabilidad de Contención (POC), que representa la probabilidad de que la persona perdida se encuentre en diferentes puntos del área de búsqueda.

Investigaciones de referencia como las de Lin y Goodrich \cite{lin2009} y Goodrich et al. \cite{lin2010} demostraron un método donde la ruta del UAV se dirige hacia regiones con la mayor probabilidad de contener el objetivo, priorizando las áreas más prometedoras primero. Las metodologías híbridas combinan las fortalezas de CPP e IPP, particionando el área operacional en subregiones de alta probabilidad basadas en un mapa de probabilidad previa, tratando cada subregión como un nodo en un problema del viajante (TSP) para determinar la secuencia óptima de tareas de cobertura \cite{vasquez2018}, así como la coordinación adaptativa de formaciones en flotas multi-UAV orientadas a la minimización del tiempo de búsqueda \cite{Multiple_Fixed-Wing_Unmanned_Aerial_Vehicles}.

\subsection{Algoritmos Baseline en SAREnv}

El framework SAREnv \cite{drones9090628} proporciona cuatro algoritmos baseline para la evaluación comparativa de estrategias de búsqueda con drones:

\begin{itemize}
	\item \textbf{Concentric Circles:} Crea una serie de trayectorias circulares concéntricas segmentadas equitativamente entre los drones disponibles. Es un planificador exhaustivo que garantiza cobertura completa dado suficiente tiempo.
	\item \textbf{Pizza Zigzag:} Divide el área circular de búsqueda en sectores tipo ``pizza'', donde cada drone cubre un sector usando un patrón zigzag (boustrophedon). También es un planificador exhaustivo.
	\item \textbf{Greedy:} Un planificador voraz que prioriza no revisitar celdas ya exploradas, dirigiéndose hacia las ubicaciones adyacentes con mayor densidad de probabilidad. Realiza exploración aleatoria si todos los espacios vecinos han sido visitados.
	\item \textbf{Random Exploration:} Exploración aleatoria a través del entorno, sirviendo como línea base inferior para comparaciones de rendimiento.
\end{itemize}

Estos algoritmos se dividen en dos categorías: planificadores exhaustivos (Concentric Circles y Pizza Zigzag) diseñados para cubrir sistemáticamente todo el área de búsqueda, y planificadores no exhaustivos (Greedy y Random) que no garantizan cobertura completa pero pueden ser más eficientes en escenarios con restricciones de tiempo.

\section{Técnicas de Búsqueda}

Las técnicas de búsqueda para drones en SAR abarcan desde métodos clásicos heurísticos hasta enfoques avanzados de aprendizaje automático y teoría de juegos.

\subsection{Métodos Heurísticos}

Los algoritmos heurísticos son simples de ejecutar pero tienen gran sobrecarga computacional cuando se usan a gran escala. Se dividen en algoritmos de planificación de cobertura y algoritmos de planificación de rutas óptimas \cite{tan2021}.

Ejemplos de algoritmos heurísticos utilizados en planificación multi-UAV incluyen A* (A-star), el algoritmo voraz (Greedy), y el algoritmo de planificación de rutas óptimas (OPP). Los algoritmos heurísticos híbridos como Elastic Band combinado con A* utilizan bandas elásticas para la planificación general de rutas y A* para suavizar la ruta inicial mediante planificación local.

\subsection{Métodos Metaheurísticos y Bioinspirados}

Los métodos metaheurísticos inspirados en la naturaleza han demostrado una elevada eficacia en la planificación de trayectorias complejas para UAVs en búsqueda y rescate, al ser capaces de explorar paisajes de búsqueda no lineales y multimodales sin quedar atrapados en óptimos locales \cite{alqudsi2025, guo2025}. Entre las técnicas bioinspiradas y basadas en física más consolidadas destacan:

\begin{itemize}
	\item \textbf{Ant Colony Optimization (ACO):} 
	Inspirado en la comunicación estigmérgica de las colonias de hormigas mediante rastros de feromona. En problemas de búsqueda en tiempo mínimo (MTS), las hormigas artificiales construyen rutas de forma probabilística combinando la concentración de feromona acumulada con una función heurística de visibilidad guiada por la matriz de probabilidad de la víctima, optimizando trayectorias continuas y discretas \cite{carabaza-ACOr, carabaza_ingles}.
	
	\item \textbf{Artificial Bee Colony (ABC):} 
	Propuesto por Karaboga y Basturk \cite{karaboga2007}, simula el comportamiento de forrajeo colectivo de un enjambre de abejas melíferas. El algoritmo organiza la población en tres roles coordinados: \textit{abejas empleadas}, que explotan y perturban localmente las fuentes de alimento (trayectorias candidatas); \textit{abejas observadoras}, que seleccionan probabilísticamente las mejores fuentes mediante mecanismos de ruleta proporcionales a la acumulación de probabilidad SAR (\textit{fitness}); y \textit{abejas exploradoras}, que abandonan soluciones estancadas tras superar un límite prefijado para muestrear aleatoriamente nuevas regiones del mapa, otorgando una gran resistencia ante mínimos locales.
	
	\item \textbf{Black Hole Algorithm (BHA):} 
	Metaheurística de inspiración física formulada por Hatamlou \cite{hatamlou2012} que emula la dinámica gravitacional de los agujeros negros. La mejor solución de la población se define como el agujero negro central, ejerciendo una atracción gravitatoria sobre el resto de soluciones candidatas (estrellas o trayectorias de vuelo). Aquellas estrellas que cruzan el radio del horizonte de sucesos son absorbidas y reubicadas estocásticamente en el espacio operacional, garantizando un equilibrio eficiente entre convergencia acelerada y exploración espacial en misiones críticas de rescate.
	
	\item \textbf{Particle Swarm Optimization (PSO) y Algoritmos Genéticos (GA):} 
	Paradigmas clásicos de inteligencia de enjambre y computación evolutiva utilizados habitualmente para la generación de waypoints y asignación de áreas en flotas multi-UAV \cite{alqudsi2025, huang2024}.
\end{itemize}

\subsection{Aprendizaje por Refuerzo}

El aprendizaje por refuerzo (RL) y el aprendizaje profundo por refuerzo (DRL) han emergido como enfoques prominentes para la planificación de rutas multi-UAV en escenarios de rescate silvestre de tiempo crítico \cite{hayat2020, yue2025, ewers2025deep}. El proceso de toma de decisiones en RL involucra factores como agente, entorno, estado, acción, recompensa, política y función de valor.

Q-learning es una técnica destacada de RL, aunque presenta problemas de sobreestimación que pueden afectar la eficiencia del sistema. Para abordar esto, se han desarrollado variantes como Deep Q-Network (DQN), Double DQN, Dueling Deep Q-Learning y C51 Dueling Deep Q-Network.

El algoritmo DDPG (Deep Deterministic Policy Gradient) muestra mejor rendimiento abordando problemas de consumo energético y convergencia comparado con el mecanismo de retrazado. MADDPG (Multi-Agent Deep Deterministic Policy Gradient) puede planificar rutas para enjambres de UAVs en entornos complejos \cite{jeon2024}.

Investigaciones consolidadas han propuesto enfoques avanzados de navegación autónoma con evasión de obstáculos basados en aprendizaje por refuerzo profundo explicable para UAVs \cite{wang2025}, equilibrando explotación y exploración en entornos continuos.

\subsection{Exploración Ergódica}

La exploración ergódica busca asignar el esfuerzo de búsqueda proporcionalmente a la distribución de probabilidad subyacente, asegurando que el tiempo dedicado a cada región sea proporcional a la densidad espacial de personas perdidas \cite{miller2016, ivic2019}.

Miller et al. \cite{miller2016} presentaron el concepto de exploración ergódica de información distribuida, donde los agentes generan trayectorias que pasan tiempo en áreas proporcional a la información esperada en esas regiones. Este enfoque ha sido extendido a múltiples agentes y ha demostrado ser efectivo para la búsqueda y rescate.

Investigaciones clave \cite{rao2023, seewald2024} han desarrollado formulaciones de optimización ergódica continua y multiagente que determinan la trayectoria óptima para cobertura considerando estrictas restricciones de batería en vehículos aéreos.

El método HEDAC (Heat equation-driven area coverage) \cite{hedac2025} es una técnica de control ergódico basada en la ecuación de calor que genera campos potenciales para dirigir enjambres en coberturas de áreas y búsqueda de personas en terrenos agrestes.

\subsection{POMDP y Procesos de Decisión}

Los Procesos de Decisión de Markov Parcialmente Observables (POMDP) proporcionan un marco matemático para modelar problemas de toma de decisiones secuenciales en entornos inciertos donde el estado del sistema no es completamente observable \cite{ragi2018}.

En SAR con UAVs, los POMDPs capturan incertidumbres en ubicaciones de objetivos, observaciones de sensores y condiciones ambientales mientras optimizan rutas. Modelan estados ambientales desconocidos, información imperfecta del sensor y la compleja interdependencia entre decisiones y observaciones futuras.

Diversos trabajos han formulado la planificación de rutas UAV-SAR mediante marcos POMDP en tiempo real, proponiendo algoritmos como Shrinking POMCP para búsqueda activa de objetivos bajo restricciones temporales \cite{zhang2024}. Por su parte, la navegación autónoma y localización colaborativa con flotas multi-UAV ha sido demostrada eficazmente para operaciones SAR en entornos con obstáculos y sin cobertura GPS \cite{sandino2020}.

Los Dec-POMDPs (Decentralized POMDPs) permiten a los UAVs tomar decisiones bajo incertidumbre con observabilidad limitada, entrenados offline usando PPO (Proximal Policy Optimization) para maximizar la suma esperada de recompensas descontadas.

\section{Métricas de Evaluación}

La evaluación de estrategias de búsqueda con drones requiere métricas cuantitativas que capturen tanto la eficiencia como la efectividad de los algoritmos de planificación.

\subsection{Probabilidad de Detección Acumulada}

La métrica de Probabilidad de Detección Acumulada ($L(\pi)$) cuantifica la probabilidad de detectar una persona perdida dada una ruta planificada $\pi$ y una distribución de probabilidad previa sobre la región de búsqueda. Se calcula como la integral de la probabilidad previa sobre las porciones del entorno visibles desde la ruta:

\begin{equation}
L(\pi) = \int_{A_\pi} P_{prior}(x) \, dx
\end{equation}

donde $A_\pi$ denota el área visible desde $\pi$, y $P_{prior}(x)$ es la probabilidad previa de que una persona perdida esté en la ubicación $x$. Esta formulación ha sido utilizada en diversos trabajos de planificación SAR donde las funciones de costo se diseñan para optimizar la coordinación multi-robot y minimizar la probabilidad acumulada de no detectar el objetivo en rescates agrestes \cite{li2020}.

\subsection{Probabilidad de Detección con Descuento Temporal}

Para incentivar el descubrimiento temprano de regiones de alta probabilidad, se introduce un factor de descuento temporal en la ganancia de probabilidad acumulada. Específicamente, para cada segmento de área $x$ observado en el tiempo $t$, la contribución al puntaje se descuenta por $\lambda^t$, donde $\lambda \in (0,1)$ controla la severidad del descuento:

\begin{equation}
I(\pi) = \int_{A_\pi} \lambda^{t(x)} P_{prior}(x) \, dx
\end{equation}

El factor de descuento $\lambda$ se calcula basado en un tiempo terminal $t_{end}$ (como el límite crítico de supervivencia). Para asegurar que el 95\% del valor acumulado sobre un horizonte infinito se capture en $t_{end}$, el factor se define como:
\begin{equation}
\lambda = (0.05)^{\frac{1}{t_{end}}}
\end{equation}
Este cálculo de $\lambda$ es especialmente importante en dominios como operaciones SAR donde la certeza de la ubicación de una persona perdida exhibe decaimiento temporal.

\subsection{Puntaje de Descubrimiento de Personas Perdidas}

El puntaje de descubrimiento de personas perdidas mide el número de personas perdidas sintéticas descubiertas durante una ruta $\pi$. En cada caso de prueba, se muestrea un número fijo de personas perdidas de una distribución predefinida, y $\pi$ se evalúa basándose en el número de personas localizadas exitosamente:

\begin{equation}
D(\pi) = \sum_{i=1}^{N} \mathbb{I}(\pi, v_i)
\end{equation}

donde la función indicadora $\mathbb{I}(\pi, v_i)$ determina si la víctima $v_i$ es localizada en la trayectoria $\pi$:
\begin{equation}
\mathbb{I}(\pi, v_i) = \begin{cases} 
1, & \text{si la víctima } v_i \text{ es descubierta durante la ruta } \pi \\
0, & \text{en caso contrario}
\end{cases}
\end{equation}

Los conteos brutos de personas encontradas han sido ampliamente utilizados en entornos de benchmarking como RoboCup Rescue \cite{kohlbrecher2015}. Formulaciones más recientes incluyen variantes ponderadas por tiempo donde las detecciones tempranas producen puntajes más altos, optimizando tanto el conteo de descubrimientos como el tiempo de descubrimiento mediante esquemas de planificación heurística y multi-objetivo \cite{trojanowski2024}.

\subsection{Probabilidad de Éxito (POS)}

Según el manual IAMSAR, la Probabilidad de Éxito (POS) es el producto de la Probabilidad de Contención (POC) y la Probabilidad de Detección (POD):

\begin{equation}
POS = POC \times POD
\end{equation}

La POC es la probabilidad de que el objetivo esté contenido en una cierta celda, mientras que el POD es la probabilidad de que el objetivo sea detectado asumiendo que está dentro de la celda cubierta por el UAV (en el campo de visión de su sensor). Cuando una celda es visitada múltiples veces, los valores de POC y POS pueden cambiar según la fórmula de actualización bayesiana \cite{koopman1980, stone2006}.

\section{Modelos de Comportamiento de Personas Perdidas}

Los modelos de comportamiento de personas perdidas son fundamentales para generar mapas de probabilidad informados que guíen la planificación de rutas de búsqueda.

\subsection{ISRID y Lost Person Behavior}

El International Search \& Rescue Incident Database (ISRID) fue desarrollado por el Dr. Robert Koester como parte de una subvención del USDA que permitió la recolección y análisis de estadísticas SAR de todo el mundo. Hasta la fecha, se han recolectado más de 150,000 incidentes SAR.

El libro ``Lost Person Behavior'' \cite{koester2008lost} utiliza clusters de patrones de comportamiento para predecir dónde podría encontrarse una persona perdida. Incluye categorías de sujetos, perfiles conductuales, estadísticas actualizadas, tareas iniciales sugeridas y preguntas de investigación especializadas.

\subsection{Modelos de Desplazamiento y Dispersión Espacial}

Para transformar las estadísticas empíricas de incidentes SAR en distribuciones espaciales de probabilidad $P_{prior}(x,y)$, la metodología formulada por Robert Koester \cite{koester2008lost}, respaldada operativamente por manuales de salvamento terrestre \cite{romero2018manual}, define cinco modelos cuantitativos y geométricos de comportamiento:

\begin{enumerate}
	\item \textbf{Modelo de Desplazamiento Radial (\textit{Ring Model}):} 
	Es el modelo primario de contención probabilística espacial. Modela la distancia euclídea radial $r = \|\mathbf{x} - \mathbf{x}_{IPP}\|$ recorrida por el sujeto desde el Punto de Planificación Inicial (IPP, \textit{Initial Planning Point}) o último punto de contacto conocido (LKP). La distribución empírica de distancias no se comporta como una gaussiana simétrica convencional, sino que presenta una acusada asimetría positiva con colas pesadas, ajustándose rigurosamente a una función de densidad log-normal:
	\begin{equation}
	f(r; \mu, \sigma) = \frac{1}{r \sigma \sqrt{2\pi}} \exp \left( -\frac{(\ln r - \mu)^2}{2\sigma^2} \right), \quad r > 0
	\end{equation}
	A partir de dicha formulación, el espacio operacional se estratifica en regiones circulares delimitadas por los radios cuartiles de dispersión estadística correspondientes a los percentiles del 25\%, 50\%, 75\% y 95\% de contención acumulada.
	
	\item \textbf{Modelo de Desviación respecto a Vías (\textit{Track Offset Model}):} 
	Cuantifica la probabilidad de que la persona extraviada permanezca a una distancia perpendicular $d_\perp(\mathbf{x}, \mathcal{L})$ respecto a una característica lineal $\mathcal{L}$ (caminos, pistas forestales, senderos o arroyos) extraída de cartografía vectorial. Las estadísticas de la base de datos ISRID evidencian que, por ejemplo, más del 50\% de las personas con demencia no halladas sobre la vía se encuentran a menos de 15 metros de su trazado, convirtiendo las bandas adyacentes a las rutas en sectores de alta probabilidad condicional.
	
	\item \textbf{Modelo de Dispersión Angular (\textit{Dispersion Model}):} 
	Modela la desviación angular del sujeto respecto al vector de dirección de viaje previsto o estimado. Individuos con desorientación severa o deterioro cognitivo tienden a exhibir ángulos de dispersión notablemente acotados debido a patrones de desplazamiento en línea recta sin cambio de rumbo, mientras que sujetos recreativos suelen presentar derivas angulares dispersas al intentar sortear obstáculos naturales.
	
	\item \textbf{Modelo de Elevación y Gradiente Topográfico:} 
	Describe la propensión relativa del desaparecido a ascender o descender en cota respecto a la elevación del IPP. Históricamente, el patrón dominante consistía en descender hacia valles siguiendo la ley de mínima resistencia gravitatoria; no obstante, el análisis contemporáneo de ISRID refleja una marcada inversión de tendencia hacia cotas altas en excursionistas para obtener visibilidad o cobertura de telefonía móvil.
	
	\item \textbf{Modelo de Cuenca Hidrográfica (\textit{Watershed Model}):} 
	Restringe físicamente la dispersión geográfica considerando las divisorias topográficas y cuencas fluviales. Empíricamente, más del 50\% de las víctimas son localizadas dentro de la misma cuenca de drenaje en la que se situaba el IPP, actuando los relieves pronunciados como barreras efectivas de contención.
\end{enumerate}

\begin{table}[H]
	\ttabbox[\textwidth]
	{\caption{PARÁMETROS ESTADÍSTICOS DE DISPERSIÓN RADIAL Y DESPLAZAMIENTO SEGÚN ISRID \cite{koester2008lost}}}
	{\label{tab:isrid_perfiles_cap2}
	\resizebox{\textwidth}{!}{%
	\begin{tabular}{|l|c|c|c|c|l|}
		\hline
		\textbf{Categoría / Perfil} & \textbf{25\% (km)} & \textbf{50\% (km)} & \textbf{75\% (km)} & \textbf{95\% (km)} & \textbf{Comportamiento de Vía (\textit{Track Offset})} \\
		\hline
		Senderista / Excursionista (\textit{Hiker}) & 0.6 & 1.8 & 3.2 & 9.9 & Alta retención en sendas ($\le 10$ m del camino) \\
		Demencia (\textit{Dementia}) & 0.3 & 1.0 & 2.4 & 12.8 & Cruza caminos al azar ($\le 15$ m en margen) \\
		Desorden del Espectro Autista (\textit{Autism}) & 0.6 & 1.6 & 3.7 & 15.2 & Atracción por cuerpos de agua y estructuras \\
		\hline
	\end{tabular}%
	}}
\end{table}

\subsection{Generación de Mapas de Probabilidad}

SAREnv \cite{drones9090628} utiliza un modelo probabilístico basado en estas estadísticas de comportamiento de personas perdidas. El proceso comienza extrayendo características geoespaciales de OpenStreetMap (estructuras, caminos, drenajes, agua, vegetación) y asigna probabilidades según tablas estadísticas de manuales.

La distribución log-normal proporciona el mejor ajuste a los datos de distancia (RMSE = 1.76 km), superando a las distribuciones normal (6.83 km) y gamma (2.02 km). Las probabilidades finales se normalizan para producir un mapa de Probabilidad de Área (POA) válido \cite{doherty2014}.

\section{Mapas Vectorizados y Datos Geoespaciales}

Los mapas vectorizados y los datos geoespaciales son la base sobre la cual se construyen los modelos de probabilidad y se planifican las rutas de búsqueda.

\subsection{OpenStreetMap}

OpenStreetMap (OSM) es la fuente principal de datos geoespaciales para aplicaciones de drones. Los datos de OSM se utilizan para extraer características relevantes para SAR como caminos, edificios, cuerpos de agua, vegetación y límites de cuencas hidrográficas.

Las herramientas basadas en OSM incluyen motores de enrutamiento como GraphHopper (Java), OSRM (Open Source Routing Machine en C++), y Valhalla (por Mapbox, soporta transporte multimodal). Para geocodificación se utiliza Nominatim.

\subsection{Procesamiento y Rasterización de Datos Vectoriales}

Para su empleo en algoritmos robóticos de planificación basados en cuadrículas, las entidades vectoriales georreferenciadas (puntos, polilíneas y polígonos) deben transformarse en matrices discretas espacialmente explícitas. Esta representación en rejilla regular (\textit{grid map}) proporciona una estructura computacional uniforme, compatible con los espacios de estados matriciales de los planificadores de búsqueda.

En implementaciones preliminares de la literatura, como el diseño original de SAREnv \cite{drones9090628}, la rasterización combinaba las capas espaciales aplicando un operador de máximo ($\max$) celda a celda para evitar saturaciones de probabilidad. No obstante, como señalan los propios autores, dicho operador introduce una distorsión operativa importante: infravalora las regiones ricas en características acumulativas (\textit{feature-rich zones}), al ignorar la convergencia de múltiples capas favorables (por ejemplo, la intersección de un sendero con una edificación de refugio). Para solventar esta limitación, las metodologías avanzadas de modelado multicapa procesan cada clase topológica en matrices binarias independientes $\mathbf{M}_k$, aplicando posteriormente una agregación lineal ponderada por los pesos del perfil conductual y una renormalización global estricta, garantizando la conservación unitaria de la creencia $\sum_{i,j} P(i,j) = 1.0$.

\subsection{Sistemas de Información Geográfica (GIS) y Proyecciones en Operaciones SAR}

La integración de herramientas GIS y librerías geoespaciales abiertas (como \texttt{OSMnx}, \texttt{GeoPandas} y \texttt{Shapely}) resulta indispensable para gestionar la complejidad espacial en misiones SAR:

\begin{itemize}
	\item \textbf{Transformación de Sistemas de Coordenadas (CRS):} 
	Los datos globales de posicionamiento GPS y cartografía OSM se expresan en coordenadas angulares geodésicas (WGS84, $\text{EPSG:4326}$). Dado que las ecuaciones cinemáticas de vuelo del UAV y los sensores requieren distancias y velocidades en el sistema métrico euclídeo sin distorsión angular, es imprescindible realizar una proyección conforme a coordenadas métricas planas Universal Transverse Mercator (UTM, como $\text{EPSG:32630}$ para la Península Ibérica).
	
	\item \textbf{Indexación Espacial Jerárquica (R-Tree):} 
	Durante la fase de preparación del escenario, la comprobación exhaustiva de intersecciones punto a polígono celda a celda conlleva una complejidad computacional prohibitiva $\mathcal{O}(N \times M)$ en rejillas de gran escala. El empleo de árboles espaciales R-Tree (\texttt{sindex}) reduce la complejidad a $\mathcal{O}(\log M)$, permitiendo identificar y enmascarar instantáneamente polígonos no transitables (lagos y estructuras cerradas).
\end{itemize}

\section{Conclusiones del Estado del Arte}

El estado del arte en enrutamiento inteligente para drones en operaciones SAR presenta un panorama diverso y en constante evolución. Las principales conclusiones que se pueden extraer son:

\begin{enumerate}
	\item \textbf{Diversidad de enfoques:} Existe una amplia variedad de técnicas de planificación de rutas, desde métodos clásicos de cobertura exhaustiva (CPP) hasta enfoques informados (IPP) y algoritmos basados en aprendizaje automático.
	
	\item \textbf{Importancia de los modelos probabilísticos:} Los modelos de comportamiento de personas perdidas, como los derivados de ISRID, son fundamentales para generar mapas de probabilidad que guíen eficazmente la búsqueda.
	
	\item \textbf{Métricas estandarizadas:} El framework SAREnv proporciona métricas estandarizadas (Probabilidad de Detección Acumulada, con descuento temporal, y Puntaje de Descubrimiento) que permiten comparar objetivamente diferentes algoritmos.
	
	\item \textbf{Tendencia hacia la inteligencia artificial:} Los métodos de aprendizaje por refuerzo y aprendizaje profundo están ganando protagonismo, especialmente para coordinación multi-agente y adaptación a entornos dinámicos.
	
	\item \textbf{Brecha simulación-realidad:} Aunque existen frameworks de simulación avanzados como SAREnv, la transferencia de algoritmos a escenarios reales sigue siendo un desafío activo de investigación.
\end{enumerate}

Este estado del arte proporciona la base teórica necesaria para el desarrollo del TFM, que se centrará en el enrutamiento inteligente sobre mapas vectorizados utilizando el framework SAREnv como plataforma de evaluación.
